\subsubsection{The Fundamental Solution for the Wave Equation}\label{subsec:wave-fourier-fundamental}

The Fourier transform is obviously useful in obtaining fundamental solutions. If $L(D)$ is a constant coefficient operator, then the equation $L(D)u=\delta$ is transformed to $L(i\bm{\xi})\widehat u=(2\pi)^{-m/2}$, i.e., to a purely algebraic equation. We immediately obtain
\begin{equation}\label{eq:fourier-fundamental-solution-symbol-reciprocal}
\widehat u(\bm{\xi})=\dfrac1{(2\pi)^{m/2}L(i\bm{\xi})},
\end{equation}
the only problem is that $L(i\bm{\xi})$ may have zeros. If \eqref{eq:fourier-fundamental-solution-symbol-reciprocal} has nonintegrable singularities, we have to consider appropriate regularizations. Finally, one has to compute the inverse Fourier transform of $\widehat u(\bm{\xi})$; this step is not necessarily easy.

Similarly, the Fourier transform can be used to find fundamental solutions for initial-value problems; we shall now do so for the wave equation in $\mathbb{R}^3$. The problem
\begin{equation}\label{eq:wave-three-dimensional-fundamental-initial-problem}
G_{tt}=\Delta G,\quad G(\bm{x},0)=0,\quad G_t(\bm{x},0)=\delta(\bm{x})
\end{equation}
is Fourier transformed in the spatial variables only; i.e., we define
\begin{equation}\label{eq:wave-spatial-fourier-transform}
\widehat G(\bm{\xi},t)=(2\pi)^{-m/2}\int_{\mathbb{R}^m}e^{-i\bm{\xi}\cdot\bm{x}}G(\bm{x},t)\,\mathrm{d}\bm{x},
\end{equation}
and apply the same type of transform to \eqref{eq:wave-three-dimensional-fundamental-initial-problem}. The result is an ODE in the variable $t$,
\begin{equation}\label{eq:wave-fourier-transformed-oscillator-problem}
\widehat G_{tt}(\bm{\xi},t)=-|\bm{\xi}|^2\widehat G(\bm{\xi},t),\quad\widehat G(\bm{\xi},0)=0,\quad\widehat G_t(\bm{\xi},0)=(2\pi)^{-3/2}.
\end{equation}
With $|\bm{\xi}|=\rho$, the solution is easily obtained as
\begin{equation}\label{eq:wave-fourier-oscillator-solution}
\widehat G(\bm{\xi},t)=(2\pi)^{-3/2}\dfrac{\sin\rho t}\rho.
\end{equation}
Using Example~\ref{ex:fourier-uniform-sphere-mass} above, we find
\begin{equation}\label{eq:wave-three-dimensional-spherical-delta}
G(\bm{x},t)=\dfrac{\delta(r-t)}{4\pi t}.
\end{equation}


It can be shown that, in any odd space dimension greater than 1, the fundamental solution of the wave equation can be expressed in terms of derivatives of $\delta(r-t)$; since there is little applied interest in solving the wave equation in more than three dimensions, we shall not prove this here. It is, however, of interest to solve the wave equation in two dimensions. In even space dimensions, it is not easy to evaluate the inverse Fourier transform of $\sin\rho t/\rho$ directly; instead, one uses a trick known as the method of descent. This trick is based on the simple observation that any solution of the wave equation in two dimensions can be regarded as a solution in three dimensions, simply by taking the direct product with the constant function 1. The fundamental solution in two dimensions can therefore be obtained by convolution of \eqref{eq:wave-three-dimensional-spherical-delta} with $\delta(x)\delta(y)1(z)$. Using the definition of convolution \eqref{eq:distribution-convolution-formal-definition}, we compute


\begin{equation}\label{eq:wave-descent-convolution-pairing}
\begin{aligned}
&\left(\delta(x)\delta(y)1(z)\dfrac{\delta(r'-t)}{4\pi t},\phi(\bm{x}+\bm{x}')\right)\\
&\quad=\dfrac1{4\pi t}\int_{-\infty}^\infty\int_{r'=t}\phi(x',y',z'+z)\,\mathrm{d}S'\,\mathrm{d}z.
\end{aligned}
\end{equation}
With $\psi(x,y)$ denoting $\int_{-\infty}^\infty\phi(x,y,z)\,\mathrm{d}z$, \eqref{eq:wave-descent-convolution-pairing} simplifies to
\begin{equation}\label{eq:wave-descent-surface-pairing}
\dfrac1{4\pi t}\int_{r'=t}\psi(x',y')\,\mathrm{d}S',
\end{equation}
and evaluation of this integral yields
\begin{equation}\label{eq:wave-descent-disk-pairing}
\dfrac1{2\pi}\int_{x^2+y^2\leq t^2}\dfrac{\psi(x,y)}{\sqrt{t^2-x^2-y^2}}\,\mathrm{d}x\,\mathrm{d}y.
\end{equation}
We have thus obtained the following fundamental solution in two space dimensions:
\begin{equation}\label{eq:wave-two-dimensional-fundamental-solution}
G(\bm{x},t)=\begin{cases}1/2\pi\sqrt{t^2-r^2},&r<t\\0,&r\geq t.\end{cases}
\end{equation}


We note that the qualitative nature of the fundamental solution for the heat equation does not really change with the space dimension, but the fundamental solution of the wave equation changes dramatically. In any number of dimensions, the support of the fundamental solution for the wave equation is contained in $|\bm{x}|\leq t$, but otherwise the fundamental solutions look quite different. Whereas the fundamental solution in three dimensions is supported only on the sphere $|\bm{x}|=t$, the support of \eqref{eq:wave-two-dimensional-fundamental-solution} fills out the full circle. Television sets in Abbott\textquotesingle s \emph{Flatland}~\cite{Ab} would have to be designed quite differently from ours; in this context, see also~\cite{Mo}.
